% =============================================================================
%  Mathematical Model Documentation — CALION: Integrated Optimization for Electrified Industrial Heating
%  Ready for Overleaf / pdfLaTeX compilation
%  Packages required: amsmath, amssymb, booktabs, array, hyperref, natbib
% =============================================================================

\documentclass[a4paper,12pt]{article}

% --- Encoding & Language ---
\usepackage[T1]{fontenc}
\usepackage[utf8]{inputenc}
\usepackage[english]{babel}

% --- Math ---
\usepackage{amsmath}
\usepackage{amssymb}
\usepackage{mathtools}          % \coloneqq, \prescript, etc.

% --- Tables ---
\usepackage{booktabs}
\usepackage{array}
\usepackage{tabularx}

% --- Figures (if needed later) ---
\usepackage{graphicx}

% --- Hyperlinks ---
\usepackage[colorlinks=true, linkcolor=blue, citecolor=blue, urlcolor=blue]{hyperref}

% --- Bibliography ---
\usepackage[numbers,sort&compress]{natbib}

% --- Misc ---
\usepackage{microtype}
\usepackage{geometry}
\geometry{margin=2.5cm}

% --- Convenience commands ---
\newcommand{\COP}{\mathrm{COP}}
\newcommand{\LMTD}{\mathrm{LMTD}}
\newcommand{\CRF}{\mathrm{CRF}}

% =============================================================================
\begin{document}

\title{Mathematical Model Documentation\\[0.5em]
       \large CALION — Integrated Optimization for Electrified Industrial Heating}
\author{}
\date{\today}
\maketitle
\tableofcontents
\newpage

% =============================================================================
\section{Model Overview and Problem Statement}
\label{sec:overview}

CALION solves a \textbf{Mixed-Integer Linear Program (MILP)} for the joint investment planning
and operational optimization of electrified industrial heating systems. The formulation follows the
\emph{unit-commitment with investment decisions} paradigm~\cite{Morales2014,Savelsbergh1994}:

\begin{equation}
  \min_{x \in \mathcal{X}} \; c^{\top} x
  \quad \text{s.t.} \quad
  Ax \leq b,\quad x \in \{0,1\}^{p} \times \mathbb{R}^{q}_{\geq 0}
  \label{eq:MILP-general}
\end{equation}

where $\mathcal{X}$ encodes:
\begin{itemize}
  \item \textbf{Binary} variables: component investment decisions $y_{i} \in \{0,1\}$ and
        operating-mode selectors;
  \item \textbf{Continuous} variables: power flows, temperatures, stored energy, fuel consumption;
  \item \textbf{Equality} constraints: energy balances (bus constraints);
  \item \textbf{Inequality} constraints: capacity limits, efficiency relationships, grid constraints.
\end{itemize}

The model horizon spans $T$ timesteps of duration $\Delta t$ hours (typically hourly resolution for
annual optimization or sub-hourly for short-horizon studies).

\paragraph{Key references.} Lund et al.~\cite{Lund2014}, Connolly et al.~\cite{Connolly2014},
Henning \& Palzer~\cite{Henning2014}.

% =============================================================================
\section{Objective Function}
\label{sec:objective}

\textit{Source files:} \texttt{calion/models/cost\_calculator.py},
\texttt{calion/models/investment\_calculator.py}.

The model minimizes the total annualized system cost over the optimization horizon:

\begin{equation}
  \min \; Z = C_{\text{energy}} + C_{\text{fuel}} + C_{\text{dump}}
             + C_{\text{CO}_2} + C_{\text{demand}}
             + C_{\text{CAPEX}} + C_{\text{act}} + C_{\text{tie}}
  \label{eq:objective}
\end{equation}

% -----------------------------------------------
\subsection{Energy Cost (Grid Electricity)}

\begin{equation}
  C_{\text{energy}}
  = \Delta t \sum_{t=1}^{T}
    \Bigl[ P_{\text{buy},t} \cdot \lambda^{\text{buy}}_{t}
          - P_{\text{sell},t} \cdot \lambda^{\text{sell}}_{t} \Bigr]
  \quad [\text{EUR}]
  \label{eq:energy-cost}
\end{equation}

\noindent Buy price:
\begin{equation}
  \lambda^{\text{buy}}_{t}
  = \lambda^{\text{spot}}_{t} + \lambda^{\text{grid}} + \lambda^{\text{levy}}
  \quad [\text{EUR/MWh}]
  \label{eq:buy-price}
\end{equation}

\noindent Sell price (reflecting market haircut and floor price):
\begin{equation}
  \lambda^{\text{sell}}_{t}
  = \max\!\Bigl(
      \lambda^{\text{spot}}_{t} \cdot (1 - h_{\text{haircut}})
      - \lambda^{\text{spread}},\;
      \lambda^{\text{floor}}
    \Bigr)
  \quad [\text{EUR/MWh}]
  \label{eq:sell-price}
\end{equation}

\noindent\textit{MILP linearization of} $\max(\cdot)$: Since the sell price is a
\emph{parameter} (computed once before model build from the price timeseries),
the $\max$ is evaluated in Python at data-loading time and passed as a constant
to the solver. No auxiliary binary variable is required.

\noindent\textit{Assumption:} The haircut fraction $h_{\text{haircut}}$ represents market
transaction costs and bilateral contract discounts; typical values 2–10\,\%~\cite{Kirschen2018}.

% -----------------------------------------------
\subsection{Fuel Cost}

\begin{equation}
  C_{\text{fuel}}
  = \Delta t \sum_{t=1}^{T} \sum_{f \in \mathcal{F}} F_{f,t} \cdot \lambda_{f}
  \quad [\text{EUR}]
  \label{eq:fuel-cost}
\end{equation}

where $\mathcal{F}$ is the set of fuel types (natural gas, biomass, waste heat),
$F_{f,t}$~[MW] is fuel consumption, and $\lambda_{f}$~[EUR/MWh] is the fuel price.

% -----------------------------------------------
\subsection{Heat Dump Penalty}

\begin{equation}
  C_{\text{dump}}
  = \Delta t \sum_{t=1}^{T} Q^{\text{dump}}_{t} \cdot \lambda_{\text{dump}}
  \quad [\text{EUR}]
  \label{eq:dump-cost}
\end{equation}

The dump cost $\lambda_{\text{dump}}$ is a penalty term ensuring that curtailed heat is
minimized; it does not represent a physical cost but guides the optimizer~\cite{Pfenninger2014}.

% -----------------------------------------------
\subsection{CO\texorpdfstring{$_2$}{2} Cost}

See Section~\ref{sec:co2}.

% -----------------------------------------------
\subsection{Peak Demand Charge}

\begin{equation}
  C_{\text{demand}}
  = \lambda_{\text{demand}} \cdot \frac{T \cdot \Delta t}{8760}
    \cdot P_{\text{buy}}^{\text{peak}}
  \quad [\text{EUR}]
  \label{eq:demand-charge}
\end{equation}

where $\lambda_{\text{demand}}$~[EUR/MW/year] is the annual demand charge rate.

\noindent\textit{Assumption:} German grid tariff structure (\emph{Leistungspreis}) applied as
a constant annual rate; methodology follows~\cite{BNetzA2023}.

% -----------------------------------------------
\subsection{CAPEX and Activation Cost}

See Section~\ref{sec:investment}.

% =============================================================================
\section{Energy Bus Balance Constraints}
\label{sec:bus}

\textit{Source file:} \texttt{calion/models/constraint\_builder.py}.

% -----------------------------------------------
\subsection{Heat Bus (Copperplate / Single Node)}

At every timestep $t$, total heat generation equals heat demand plus losses:

\begin{equation}
  \sum_{i \in \mathcal{G}_{\text{th}}} Q^{\text{out}}_{i,t}
  = Q^{\text{demand}}_{t} + Q^{\text{dump}}_{t}
    + \sum_{j \in \mathcal{S}} Q^{\text{charge}}_{j,t}
    + Q^{\text{net\,loss}}_{t}
  \quad \forall\, t
  \label{eq:heat-bus}
\end{equation}

where $\mathcal{G}_{\text{th}}$ = thermal generators (boilers, heat pumps, CHP, P2H, storage
discharge) and $\mathcal{S}$ = thermal storage units.

\noindent\textit{Physical interpretation:} Conservation of energy (first law of thermodynamics)
applied at the thermal bus~\cite{Lund2014}.

% -----------------------------------------------
\subsection{Electricity Bus}

\begin{equation}
  P^{\text{buy}}_{t} + \sum_{i \in \mathcal{G}_{\text{el}}} P^{\text{el,out}}_{i,t}
  = \sum_{j \in \mathcal{C}_{\text{el}}} P^{\text{el,in}}_{j,t} + P^{\text{sell}}_{t}
  \quad \forall\, t
  \label{eq:el-bus}
\end{equation}

where $\mathcal{G}_{\text{el}}$ = CHP units (electricity generation),
$\mathcal{C}_{\text{el}}$ = electricity consumers (heat pumps, P2H).

% -----------------------------------------------
\subsection{Per-Node Heat Balance (Multi-Node Networks)}

\noindent\textbf{Producer node:}
\begin{equation}
  \sum_{i} Q^{\text{out}}_{i,t}
  = Q^{\text{demand}}_{t} + Q^{\text{dump}}_{t}
    + \sum_{j} Q^{\text{charge}}_{j,t}
    + Q^{\text{pipe\,out}}_{t} - Q^{\text{pipe\,in}}_{t}
    + Q^{\text{net\,loss}}_{t}
  \label{eq:producer-node}
\end{equation}

\noindent\textbf{Consumer node:}
\begin{equation}
  Q^{\text{pipe\,delivered}}_{t} + \sum_{i} Q^{\text{local,out}}_{i,t} + Q^{\text{slack}}_{t}
  = Q^{\text{demand}}_{t} + \sum_{j} Q^{\text{charge}}_{j,t}
  \label{eq:consumer-node}
\end{equation}

where $Q^{\text{slack}}_{t} \geq 0$~[MW] is a demand-slack variable penalized in the
objective by a large coefficient $\lambda_{\text{slack}}$~[EUR/MWh] (default: $10^6$).
This slack prevents artificial infeasibility at nodes with multiple incoming pipes when
heat delivery is bounded by pipe capacity~\cite{Pfenninger2014}:

\begin{equation}
  C_{\text{slack}} = \Delta t \sum_{t} Q^{\text{slack}}_{t} \cdot \lambda_{\text{slack}}
  \quad [\text{EUR}]
  \label{eq:slack-cost}
\end{equation}

\noindent\textit{Source file:} \texttt{calion/models/network\_manager.py}
(constraint \texttt{demand\_heat\_rule}, parameter \texttt{demand\_slack\_penalty\_eur\_per\_mwh}).

These balance equations follow the nodal energy formulation, analogous to Kirchhoff's
current law applied to energy networks~\cite{Vesterlund2017,Lund2016b}.

% =============================================================================
\section{Heat Pump Model}
\label{sec:hp}

\textit{Source file:} \texttt{calion/models/blocks/heat\_pump.py}.

% -----------------------------------------------
\subsection{Heat Output Balance}

The total heat output $Q_{t}$ from a heat pump is the sum of contributions from the
waste-heat recovery (WRG) source and the default (ambient or other) source:

\begin{equation}
  Q_{t} = Q^{\text{WRG}}_{t} + Q^{\text{def}}_{t}
  \quad \forall\, t
  \tag{HP-1}
  \label{eq:hp1}
\end{equation}

% -----------------------------------------------
\subsection{Electrical Power Consumption}

\begin{equation}
  P^{\text{el}}_{t}
  = \frac{Q^{\text{WRG}}_{t}}{\COP_{t}}
    + \frac{Q^{\text{def}}_{t}}{\COP_{\text{default}}}
  \quad \forall\, t
  \tag{HP-2}
  \label{eq:hp2}
\end{equation}

\noindent\textit{Derivation:} The definition of the Coefficient of Performance is
\[
  \COP = \frac{Q_{\text{th}}}{P_{\text{el}}}
  \implies
  P_{\text{el}} = \frac{Q_{\text{th}}}{\COP}.
\]
With two heat sources at potentially different COPs, Eq.~\eqref{eq:hp2} follows
directly~\cite{Viessmann2021,Berghmans2012}.

% -----------------------------------------------
\subsection{Capacity Constraint}

\begin{equation}
  Q_{t} \leq \hat{Q} \cdot u_{t}
  \quad \forall\, t
  \tag{HP-3}
  \label{eq:hp3}
\end{equation}

where $\hat{Q}$~[MW] is the installed thermal capacity and $u_{t} \in \{0,1\}$ is the
operating binary variable.

% -----------------------------------------------
\subsection{Minimum Load Constraint}

\begin{equation}
  Q_{t} \geq \phi_{\min} \cdot \hat{Q} \cdot u_{t}
  \quad \forall\, t
  \tag{HP-4}
  \label{eq:hp4}
\end{equation}

where $\phi_{\min} \in [0,1]$ is the minimum load fraction. This models the technical
minimum operation point of the heat pump compressor~\cite{Wolf2014}.

% -----------------------------------------------
\subsection{Investment Constraints (Capacity Expansion)}

\begin{align}
  \hat{Q} &\leq \hat{Q}_{\max} \cdot y
  \tag{HP-5}
  \label{eq:hp5} \\
  \hat{Q} &\geq \hat{Q}_{\min} \cdot y
  \tag{HP-6}
  \label{eq:hp6} \\
  u_{t}   &\leq y \quad \forall\, t
  \tag{HP-7}
  \label{eq:hp7}
\end{align}

where $y \in \{0,1\}$ is the binary investment decision, and
$\hat{Q}_{\min}, \hat{Q}_{\max}$ are the minimum and maximum investable
capacities~\cite{Morales2014,Morvaj2016}.

\noindent\textit{Formulation type:} Standard big-M disjunctive constraints; see
Raman \& Grossmann~\cite{Raman1994}.

% =============================================================================
\section{COP Calculation}
\label{sec:cop}

\textit{Source file:} \texttt{calion/models/cop\_calculator.py}.

% -----------------------------------------------
\subsection{Analytical Thermodynamic COP Model}

The analytical COP model is based on the semi-empirical formula for electrically-driven
compression heat pumps, derived from second-law thermodynamic
analysis~\cite{Staffell2012,Ruhnau2019}:

\begin{align}
  \COP
  &= \underbrace{%
       \frac{L_s}{L_s - L_{\text{src}}}
     }_{A}
     \cdot
     \underbrace{%
       \frac{%
         1 + \bigl(\delta_s + \Delta T_{\text{pp}}\bigr) / L_s
       }{%
         1 + \bigl(\delta_s + \tfrac{1}{2}(T_{\text{src,in}} - T_{\text{src,out}})
             + 2\Delta T_{\text{pp}}\bigr)
             / (L_s - L_{\text{src}})
       }
     }_{B}
     \cdot \eta \cdot (1 - q_{ww})
     + (1 - \eta) - F_{Q}
  \tag{COP-1}
  \label{eq:cop1}
\end{align}

\noindent where:
\begin{itemize}
  \item $L_s$ = LMTD of heat sink [K]
  \item $L_{\text{src}}$ = LMTD of heat source [K]
  \item $\eta$ = Carnot efficiency factor (typical: 0.75)~[--]
  \item $q_{ww}$ = auxiliary loss factor (compressor, fan, pump)
  \item $F_Q$ = fixed loss factor (typical: 0.10)~[--]
  \item $\Delta T_{\text{pp}}$ = pinch-point temperature approach [K]
  \item $\delta_s$ = mean temperature approach parameter at sink
\end{itemize}

\noindent\textit{Scientific basis:} The formula structure follows the approach of
Staffell et al.~\cite{Staffell2012} and Ruhnau et al.~\cite{Ruhnau2019}, which
parameterize real heat pump performance as a fraction of the ideal Carnot COP,
corrected for internal temperature approaches.

\paragraph{Carnot COP reference.}

\begin{equation}
  \COP_{\text{Carnot}}
  = \frac{T_{\text{sink}}}{T_{\text{sink}} - T_{\text{source}}}
  \quad [K/K]
  \label{eq:cop-carnot}
\end{equation}

where temperatures are in Kelvin~\cite{Cengel2018}.

% -----------------------------------------------
\subsection{Log-Mean Temperature Difference (LMTD)}

\begin{equation}
  \LMTD(T_h, T_c)
  = \frac{T_h - T_c}{\ln(T_h / T_c)}
  \tag{COP-2}
  \label{eq:lmtd}
\end{equation}

\noindent Limit (L'H\^{o}pital): when $T_h \approx T_c$:
\[
  \LMTD \approx \frac{T_h + T_c}{2} \quad \text{(arithmetic mean)}
\]

\noindent\textit{Standard reference:} Incropera et al.~\cite{Incropera2017}, Chapter~11.

% -----------------------------------------------
\subsection{COP Lookup-Table Model (2D Bilinear Interpolation)}

When manufacturer data is available, COP is computed via bilinear interpolation over a
grid of source and sink temperatures:

\begin{equation}
  \COP(T_{\text{src}}, T_{\text{sink}})
  = \sum_{i,j} w_{ij} \cdot \COP_{ij}
  \label{eq:cop-lookup}
\end{equation}

where $w_{ij}$ are bilinear weights from the 2D table. Clamping:
\[
  \COP \in [\COP_{\min},\; \COP_{\max}] = [0.5,\; 15.0]
\]

\noindent\textit{Reference:} EN~14825:2022 standard for heat pump performance
testing~\cite{EN14825}.

% =============================================================================
\section{Thermal Storage Model}
\label{sec:storage}

\textit{Source file:} \texttt{calion/models/blocks/storage.py}.

% -----------------------------------------------
\subsection{Energy Balance (State-of-Charge Dynamics)}

The discrete-time energy balance of a thermal storage unit follows:

\begin{equation}
  E_{t}
  = E_{t-1} \cdot \lambda^{\Delta t}
    + \eta_{c} \cdot Q^{c}_{t} \cdot \Delta t
    - \frac{Q^{d}_{t} \cdot \Delta t}{\eta_{d}}
  \quad \forall\, t
  \tag{ST-1}
  \label{eq:storage-soc}
\end{equation}

\noindent where:
\begin{itemize}
  \item $E_{t}$ [MWh] = stored energy (state-of-charge) at timestep $t$
  \item $\lambda$ = self-discharge factor per hour [--], $\lambda = 1 - \dot{\lambda}_{\text{loss}}$
  \item $\lambda^{\Delta t}$ = self-discharge over one timestep (accounts for variable $\Delta t$)
  \item $\eta_{c}$ = charge efficiency [--]
  \item $\eta_{d}$ = discharge efficiency [--]
  \item $Q^{c}_{t}$ [MW] = charge power (heat input)
  \item $Q^{d}_{t}$ [MW] = discharge power (heat output)
  \item $\Delta t$ [h] = timestep duration
\end{itemize}

\noindent\textit{Derivation:} Standard energy storage model with multiplicative
self-discharge; see Quoilin et al.~\cite{Quoilin2017} and Lund et al.~\cite{Lund2007}.
The model uses $\lambda^{\Delta t}$ rather than $\lambda \cdot \Delta t$ to correctly
represent exponential decay for variable timestep lengths.

% -----------------------------------------------
\subsection{Energy Capacity Constraints}

\begin{align}
  E_{t} &\leq \hat{E} \cdot \xi^{\text{active}}_{t}    && \forall\, t \tag{ST-2} \label{eq:st2} \\
  E_{t} &\geq E_{\min} \cdot \xi^{\text{active}}_{t}   && \forall\, t \tag{ST-3} \label{eq:st3}
\end{align}

where $\hat{E}$~[MWh] is the energy capacity (decision variable in expansion mode) and
$\xi^{\text{active}}_{t} \in \{0,1\}$ is the storage activation binary.
The upper bound $E_t \leq \hat{E} \cdot \xi^{\text{active}}_t$ additionally enforces
$E_t \leq \hat{E}_{\max} \cdot \xi^{\text{active}}_t$ to handle the fixed-capacity case.

\noindent\textit{Source file:} \texttt{calion/models/blocks/storage.py} (constraints
\texttt{ecap\_hi}, \texttt{emax\_hi}, \texttt{ecap\_lo}).

% -----------------------------------------------
\subsection{Power Capacity Constraints}

\begin{align}
  Q^{c}_{t} &\leq \hat{P} \cdot \xi^{c}_{t} && \forall\, t \tag{ST-4} \label{eq:st4} \\
  Q^{d}_{t} &\leq \hat{P} \cdot \xi^{d}_{t} && \forall\, t \tag{ST-5} \label{eq:st5}
\end{align}

where $\hat{P}$~[MW] is the power capacity and $\xi^{c}_{t}, \xi^{d}_{t} \in \{0,1\}$
are binary charge/discharge mode indicators.

% -----------------------------------------------
\subsection{Simultaneous Charge/Discharge Prevention}

\begin{equation}
  \xi^{c}_{t} + \xi^{d}_{t} \leq 1
  \quad \forall\, t
  \tag{ST-6}
  \label{eq:st6}
\end{equation}

\noindent\textit{Physical justification:} Thermal storage cannot simultaneously absorb and
release heat (ignoring stratification effects). See Lund et al.~\cite{Lund2007}.

% -----------------------------------------------
\subsection{Power-Energy Coupling (Optional)}

\begin{equation}
  \hat{P} \leq \kappa \cdot \hat{E}
  \tag{ST-7}
  \label{eq:st7}
\end{equation}

where $\kappa$~[MW/MWh = h$^{-1}$] is the C-rate coupling factor.

\noindent\textit{Reference:} Common for sensible heat thermal energy storage (e.g., water
tanks); typical $\kappa$ values 0.1–1.0~\cite{Haller2019}.

% -----------------------------------------------
\subsection{Initial State-of-Charge Bound}

\begin{equation}
  E_{0} \leq \hat{E}
  \tag{ST-8}
  \label{eq:st8}
\end{equation}

The initial state-of-charge $E_0$ is a user-defined parameter that must not exceed the
installed energy capacity.

\noindent\textit{Source file:} \texttt{calion/models/blocks/storage.py}
(constraint \texttt{soc0\_cap}).

% -----------------------------------------------
\subsection{Cyclic Terminal Constraint}

For planning studies with periodic (e.g.\ annual or monthly) horizons, the storage must
return to its initial state at the end of the horizon:

\begin{equation}
  E_{T} = E_{0}
  \tag{ST-9}
  \label{eq:st9}
\end{equation}

This \emph{cyclic} boundary condition prevents the optimizer from artificially draining
storage at the end of the horizon~\cite{Lund2007}.
Alternatively, softer policies can be applied:
\begin{align*}
  \text{``equal'':} & \quad E_{T} = E_{0} \\
  \text{``geq'':}   & \quad E_{T} \geq E_{0} \\
  \text{``free'':}  & \quad \text{no terminal constraint}
\end{align*}

In all Memmingen runs the \texttt{equal} policy is used.

\noindent\textit{Source file:} \texttt{calion/models/state\_constraints.py}
(terminal constraint, applied at system-builder level).

% =============================================================================
\section{Power-to-Heat (P2H) Model}
\label{sec:p2h}

\textit{Source file:} \texttt{calion/models/blocks/p2h.py}.

% -----------------------------------------------
\subsection{Efficiency-Based Heat Output}

\begin{equation}
  Q_{t} = \eta_{t} \cdot P^{\text{el}}_{t}
  \quad \forall\, t
  \tag{P2H-1}
  \label{eq:p2h1}
\end{equation}

where $\eta_{t}$~[--] is the electrical-to-thermal conversion efficiency ($\leq 1$ for
resistive heating, $> 1$ not applicable --- unlike heat pumps).

\noindent\textit{Device types covered:} Electric boiler (immersion heater), electrode boiler,
infrared heater. Typical $\eta$ = 0.95–0.99~\cite{Lund2015}.

% -----------------------------------------------
\subsection{Capacity and Minimum Load}

\begin{align}
  Q_{t} &\leq \hat{Q} \cdot u_{t}              && \forall\, t \tag{P2H-2} \label{eq:p2h2} \\
  Q_{t} &\geq \phi_{\min} \cdot \hat{Q} \cdot u_{t} && \forall\, t \tag{P2H-3} \label{eq:p2h3}
\end{align}

% =============================================================================
\section{Thermal Generator and CHP Model}
\label{sec:gen}

\textit{Source file:} \texttt{calion/models/blocks/thermal\_gen.py}.

% -----------------------------------------------
\subsection{Boiler: Thermal Output from Fuel}

\begin{equation}
  Q^{\text{th}}_{t} = \eta^{\text{th}} \cdot F_{t}
  \quad \forall\, t
  \tag{GEN-1}
  \label{eq:gen1}
\end{equation}

where $\eta^{\text{th}}$ is the net thermal efficiency (lower heating value, LHV basis) and
$F_{t}$~[MW] is fuel input power.

\noindent\textit{Typical values:} Gas condensing boiler $\eta^{\text{th}}$ = 0.90–1.05
(LHV), biomass boiler 0.80–0.90~\cite{Eurostat2023}.

% -----------------------------------------------
\subsection{CHP: Combined Outputs}

For a combined heat and power unit, both heat and electricity are produced from the same
fuel input~\cite{Morales2014}:

\begin{align}
  Q^{\text{th}}_{t} &= \eta^{\text{th}} \cdot F_{t} && \forall\, t \tag{GEN-2} \label{eq:gen2} \\
  P^{\text{el}}_{t} &= \eta^{\text{el}} \cdot F_{t} && \forall\, t \tag{GEN-3} \label{eq:gen3}
\end{align}

The total efficiency is:
\begin{equation}
  \eta^{\text{total}} = \eta^{\text{th}} + \eta^{\text{el}}
  \label{eq:total-eff}
\end{equation}

\noindent\textit{Assumption:} Fixed extraction CHP model (backpressure turbine). More
sophisticated CHP models with a feasible operating region (FOR) or iso-fuel lines are not
currently implemented; see Morales et al.~\cite{Morales2014} for the general formulation.

% -----------------------------------------------
\subsection{Capacity Constraint}

\begin{equation}
  Q^{\text{th}}_{t} \leq \hat{Q}^{\text{th}}
  \quad \forall\, t
  \tag{GEN-4}
  \label{eq:gen4}
\end{equation}

% =============================================================================
\section{Pipe Heat Loss and Network Physics}
\label{sec:pipe}

\textit{Source files:} \texttt{calion/models/network\_physics.py},
\texttt{calion/models/blocks/pipe\_pair.py}.

% -----------------------------------------------
\subsection{Steady-State Pipe Heat Loss}

The heat loss per unit length of a buried insulated pipe (steady-state model):
\[
  \dot{q}_{\text{loss}} = U \cdot (T_{\text{fluid}} - T_{\text{ground}}) \quad [\text{W/m}]
\]

Total heat loss for a pipe of length $L$:
\begin{equation}
  Q_{\text{loss}}
  = U \cdot L \cdot (T_{\text{fluid}} - T_{\text{ground}}) \times 10^{-6}
  \quad [\text{MW}]
  \tag{PH-1}
  \label{eq:ph1}
\end{equation}

where $U$~[W/(m$\cdot$K)] is the linear heat transfer coefficient of the insulated pipe
(combined insulation + soil conductance).

\noindent\textit{Scientific reference:} EN~13941-1:2019~\cite{EN13941}, Frederiksen \&
Werner~\cite{Frederiksen2013} \S4.3.

\noindent\textbf{Supply and return separately:}
\begin{align*}
  Q^{\text{loss}}_{\text{supply}}
    &= U_s \cdot L \cdot (T_s - T_g) \times 10^{-6} \\
  Q^{\text{loss}}_{\text{return}}
    &= U_r \cdot L \cdot (T_r - T_g) \times 10^{-6}
\end{align*}

% -----------------------------------------------
\subsection{Outlet Temperature After Heat Loss}

Given inlet temperature $T_{\text{in}}$, pipe geometry, and mass flow $\dot{m}$:
\[
  Q_{\text{loss}} = U \cdot L \cdot (T_{\text{in}} - T_g) \times 10^{-3} \quad [\text{kW}]
  \qquad
  \Delta T = \frac{Q_{\text{loss}}}{\dot{m} \cdot c_p} \quad [\text{K}]
  \qquad
  T_{\text{out}} = T_{\text{in}} - \Delta T
  \tag{PH-2}
  \label{eq:ph2}
\]

or in closed form:
\begin{equation}
  T_{\text{out}}
  = T_g + (T_{\text{in}} - T_g)
    \cdot \exp\!\left(-\frac{U \cdot L}{\dot{m} \cdot c_p \cdot 1000}\right)
  \approx
  T_{\text{in}} - \frac{U \cdot L \cdot (T_{\text{in}} - T_g)}{\dot{m} \cdot c_p \cdot 1000}
  \label{eq:ph2-closed}
\end{equation}

\noindent\textit{Note:} The exponential form is the exact solution for plug-flow with linear
heat loss; the linear approximation holds when $UL \ll \dot{m} c_p$~\cite{Benonysson1995}.

\noindent\textit{Reference:} Benonysson et al.~\cite{Benonysson1995}, Frederiksen \&
Werner~\cite{Frederiksen2013} \S4.5.

% -----------------------------------------------
\subsection{Mass Flow from Heat Delivery}

The required mass flow $\dot{m}$ to deliver heat power $Q$~[kW] with supply/return
temperature difference $\Delta T$:

\begin{equation}
  \dot{m} = \frac{Q}{c_p \cdot \Delta T}
  \quad [\text{kg/s}]
  \tag{PH-3}
  \label{eq:ph3}
\end{equation}

where $c_p = 4.186$~kJ/(kg$\cdot$K) for water.

\noindent\textit{Reference:} Standard heat transport equation~\cite{Frederiksen2013}.

% -----------------------------------------------
\subsection{Flow Velocity}

\begin{equation}
  v = \frac{\dot{m}}{\rho \cdot A}
  \quad [\text{m/s}],
  \qquad
  A = \pi (d/2)^{2} \quad [\text{m}^{2}]
  \tag{PH-4}
  \label{eq:ph4}
\end{equation}

where $\rho \approx 975$~kg/m$^3$ (hot water at $\sim$70\,$^{\circ}$C) and $d$ is inner
pipe diameter [m].

% -----------------------------------------------
\subsection{MILP Linear Heat Delivery}

In the MILP formulation, temperature is held at nominal values, linearizing the
temperature-dependent heat delivery:

\begin{equation}
  Q^{\text{delivered}}_{t} \cdot 1000
  = \dot{m}_{t} \cdot c_p \cdot \Delta T_{\text{nom}}
  \tag{PH-5}
  \label{eq:ph5}
\end{equation}

\noindent\textit{Assumption:} Fixed supply/return temperature difference
($\Delta T$ assumption); standard in MILP district heating
models~\cite{Vesterlund2017,Mesfun2015}.

% -----------------------------------------------
\subsection{Heating Curve (Heizkurve)}

Weather-compensated supply temperature control:

\begin{equation}
  T_s(T_{\text{out}})
  = T_{s,\min}
    + (T_{s,\max} - T_{s,\min})
      \cdot \frac{T_{\text{out,design}} - T_{\text{out}}}{T_{\text{out,design}} - T_{\text{out,min}}}
  \tag{PH-6}
  \label{eq:heating-curve}
\end{equation}

Linear equivalent: $T_s = \alpha + \beta \cdot T_{\text{out}}$, where:
\[
  \beta = -\frac{T_{s,\max} - T_{s,\min}}{T_{\text{out,design}} - T_{\text{out,min}}}
\]

\noindent\textit{Reference:} VDI~2067~\cite{VDI2067}, Frederiksen \&
Werner~\cite{Frederiksen2013} \S6.2.

% -----------------------------------------------
\subsection{Transport Delay (SOS2 Piecewise-Linear Approximation)}

Flow-dependent thermal transport delay through pipes:

\begin{equation}
  \tau(t) = \frac{L \cdot \rho \cdot A}{\dot{m}(t)}
  \quad [\text{s}]
  \tag{PH-7}
  \label{eq:ph7}
\end{equation}

For the MILP model, the nonlinear function $\tau(\dot{m})$ is approximated using Special
Ordered Sets of type~2 (SOS2) with 3 breakpoints~\cite{Williams2013}:

\begin{table}[h!]
  \centering
  \caption{SOS2 breakpoints for transport delay approximation.}
  \begin{tabular}{@{}lll@{}}
    \toprule
    Bucket & Flow range & Delay \\
    \midrule
    0 (high flow)   & $[\dot{m}_{\text{mid}}, \dot{m}_{\text{max}}]$ & $\tau_1$ (shortest) \\
    1 (medium flow) & $[\dot{m}_{\text{min}}, \dot{m}_{\text{mid}}]$ & $\tau_2$ \\
    2 (low flow)    & $[0, \dot{m}_{\text{min}}]$                    & $\tau_3$ (longest) \\
    \bottomrule
  \end{tabular}
  \label{tab:sos2}
\end{table}

\noindent\textit{Reference:} SOS2 modeling~\cite{Beale1970}; transport delay in
DH~\cite{Schweiger2017}.

% =============================================================================
\section{PWL Temperature Linearization}
\label{sec:pwl}

\textit{Source file:} \texttt{calion/models/blocks/temperature\_linearization.py}.

In the MILP formulation all temperatures are pre-computed at model-build time as
functions of the normalized network load fraction.  Three methods are available;
the \texttt{pwl} method is used for all Memmingen runs.

% -----------------------------------------------
\subsection{Load Fraction}

\begin{equation}
  \lambda_{t}
  = \frac{Q^{\text{demand}}_{t}}{Q^{\text{peak}}}
  \in [0,\;1]
  \quad \forall\, t
  \tag{PWL-1}
  \label{eq:pwl1}
\end{equation}

where $Q^{\text{peak}}$~[MW] is the maximum network demand over the horizon.

% -----------------------------------------------
\subsection{Piecewise-Linear Temperature Profile}

Given $N+1$ breakpoints
$\{(\lambda^{(k)}, T^{(k)}_s, T^{(k)}_r)\}_{k=0}^{N}$
with $0 \leq \lambda^{(0)} < \cdots < \lambda^{(N)} \leq 1$
and $T^{(k)}_s > T^{(k)}_r$ for all $k$, the supply and return temperatures
are interpolated for each timestep $t$:

\begin{align}
  T_s(t) &= T^{(k)}_s
             + \frac{\lambda_t - \lambda^{(k)}}{\lambda^{(k+1)} - \lambda^{(k)}}
               \bigl(T^{(k+1)}_s - T^{(k)}_s\bigr)
  \tag{PWL-2a}
  \label{eq:pwl2a} \\
  T_r(t) &= T^{(k)}_r
             + \frac{\lambda_t - \lambda^{(k)}}{\lambda^{(k+1)} - \lambda^{(k)}}
               \bigl(T^{(k+1)}_r - T^{(k)}_r\bigr)
  \tag{PWL-2b}
  \label{eq:pwl2b}
\end{align}

where index $k$ is the active segment: $\lambda^{(k)} \leq \lambda_t < \lambda^{(k+1)}$.
Values are clamped to the first/last breakpoint outside the range.

% -----------------------------------------------
\subsection{Memmingen Breakpoints}

\begin{table}[h!]
  \centering
  \caption{PWL temperature profile used in all Memmingen runs (L2 and L3).}
  \begin{tabular}{@{}cccc@{}}
    \toprule
    Breakpoint $k$ & $\lambda^{(k)}$ & $T^{(k)}_s$ [$^{\circ}$C] & $T^{(k)}_r$ [$^{\circ}$C] \\
    \midrule
    0 & 0.0 & 70 & 55 \\
    1 & 0.3 & 80 & 50 \\
    2 & 0.6 & 90 & 45 \\
    3 & 1.0 & 100 & 40 \\
    \bottomrule
  \end{tabular}
  \label{tab:pwl-breakpoints}
\end{table}

\noindent\textit{Physical interpretation:} Higher network load requires higher supply
temperature (larger $\Delta T$) to deliver heat within the hydraulic capacity constraints.
The return temperature decreases with load because consumers extract more heat, cooling
the return flow further.

% -----------------------------------------------
\subsection{Pipe Loss with Load-Dependent Temperatures}

The PWL temperatures enter the pipe heat loss model (Section~\ref{sec:pipe}) at each
timestep:

\begin{align}
  Q^{\text{loss}}_{\text{supply},t}
    &= U_s \cdot L \cdot \bigl(T_s(t) - T_g\bigr) \times 10^{-6}
    \quad [\text{MW}]
    \tag{PWL-3a}
    \label{eq:pwl3a} \\
  Q^{\text{loss}}_{\text{return},t}
    &= U_r \cdot L \cdot \bigl(T_r(t) - T_g\bigr) \times 10^{-6}
    \quad [\text{MW}]
    \tag{PWL-3b}
    \label{eq:pwl3b}
\end{align}

Since $T_s(t)$ and $T_r(t)$ are pre-computed parameters (not decision variables),
the pipe loss remains a linear expression in the MILP~\cite{Vesterlund2017}.

\noindent\textit{Alternative methods:}
\begin{itemize}
  \item \texttt{fixed}: $T_s(t) \equiv T_{s,\text{nom}}$,\; $T_r(t) \equiv T_{r,\text{nom}}$
        (constant, zero computation overhead).
  \item \texttt{global}: Single linear segment between minimum and maximum load endpoints
        (special case of PWL with $N=1$).
\end{itemize}

% =============================================================================
\section{Heat Exchanger Model}
\label{sec:hx}

\textit{Source file:} \texttt{calion/models/blocks/heat\_exchanger.py}.

% -----------------------------------------------
\subsection{Energy Balances}

\noindent Primary (hot) side:
\begin{equation}
  Q^{\text{transfer}}_{t} \cdot 1000
  = \dot{m}^{\text{prim}}_{t} \cdot c_p
    \cdot \bigl(T^{\text{prim,in}}_{t} - T^{\text{prim,out}}_{t}\bigr)
  \tag{HX-1}
  \label{eq:hx1}
\end{equation}

\noindent Secondary (cold) side:
\begin{equation}
  Q^{\text{transfer}}_{t} \cdot 1000
  = \dot{m}^{\text{sec}}_{t} \cdot c_p
    \cdot \bigl(T^{\text{sec,out}}_{t} - T^{\text{sec,in}}_{t}\bigr)
  \tag{HX-2}
  \label{eq:hx2}
\end{equation}

% -----------------------------------------------
\subsection{Effectiveness Constraint (Linearized)}

For a counter-flow heat exchanger, the effectiveness $\varepsilon$ is:
\[
  \varepsilon
  = \frac{Q_{\text{actual}}}{Q_{\text{max}}}
  = \frac{T^{\text{sec,out}} - T^{\text{sec,in}}}{T^{\text{prim,in}} - T^{\text{sec,in}}}
\]

Linearized as an upper bound:
\begin{equation}
  T^{\text{sec,out}}_{t}
  \leq T^{\text{sec,in}}_{t}
       + \varepsilon \cdot \bigl(T^{\text{prim,in}}_{t} - T^{\text{sec,in}}_{t}\bigr)
  \tag{HX-3}
  \label{eq:hx3}
\end{equation}

\noindent\textit{Reference:} Incropera et al.~\cite{Incropera2017} \S11.4; NTU-effectiveness
method.

% -----------------------------------------------
\subsection{Pinch Point Constraint (Big-M Formulation)}

\begin{equation}
  T^{\text{prim,out}}_{t}
  \geq T^{\text{sec,in}}_{t} + \Delta T_{\text{pp}}
       - M \cdot (1 - a_{t})
  \tag{HX-4}
  \label{eq:hx4}
\end{equation}

where $\Delta T_{\text{pp}}$ is the minimum approach temperature [K] and $M$ is a large
constant.

\noindent\textit{Physical meaning:} Second law of thermodynamics prevents temperature cross at
the pinch point~\cite{Incropera2017}.

% =============================================================================
\section{Electricity Market Model}
\label{sec:elmarket}

\textit{Source file:} \texttt{calion/models/constraint\_builder.py}.

% -----------------------------------------------
\subsection{Buy/Sell Exclusivity}

\begin{align}
  P^{\text{buy}}_{t}  &\leq M_{G} \cdot g_{t}         && \forall\, t \tag{EL-1} \label{eq:el1} \\
  P^{\text{sell}}_{t} &\leq M_{G} \cdot (1 - g_{t})   && \forall\, t \tag{EL-2} \label{eq:el2}
\end{align}

where $g_{t} \in \{0,1\}$ is the grid mode binary variable and $M_{G}$ is a big-M constant.

\noindent\textit{Assumption:} The system is a price-taker on the electricity spot market; no
market price feedback from dispatch decisions~\cite{Kirschen2018,Hirth2013}.

% -----------------------------------------------
\subsection{Connection Limits}

\begin{align}
  P^{\text{buy}}_{t}  &\leq P^{\text{import,max}} && \forall\, t \tag{EL-3} \label{eq:el3} \\
  P^{\text{sell}}_{t} &\leq P^{\text{export,max}} && \forall\, t \tag{EL-4} \label{eq:el4}
\end{align}

% -----------------------------------------------
\subsection{Peak Demand Tracking}

\begin{equation}
  P^{\text{peak}} \geq P^{\text{buy}}_{t} \quad \forall\, t
  \tag{EL-5}
  \label{eq:el5}
\end{equation}

This linear constraint correctly identifies the maximum import for demand charge calculation
without requiring auxiliary binary variables.

% =============================================================================
\section{Investment and Annuity Model}
\label{sec:investment}

\textit{Source file:} \texttt{calion/models/investment\_calculator.py}.

% -----------------------------------------------
\subsection{Annualization Factor}

For an optimization horizon of $T$ timesteps with duration $\Delta t$ hours:
\[
  f_{\text{period}} = \frac{T \cdot \Delta t}{8760} \quad [\text{yr/yr}]
\]

Annualized cost factor (equivalent annual cost, EAC):

\begin{equation}
  f_{\text{ann}} = \frac{f_{\text{period}}}{\tau_{\text{life}}}
  \tag{INV-1}
  \label{eq:inv1}
\end{equation}

where $\tau_{\text{life}}$ is the economic lifetime in years.

\noindent\textbf{Note:} The simplified annualization in Eq.~\eqref{eq:inv1} assumes a
discount rate of zero. For a non-zero discount rate $r$, the correct Capital Recovery Factor
is:

\begin{equation}
  \CRF(r, n) = \frac{r(1+r)^{n}}{(1+r)^{n} - 1}
  \tag{INV-2}
  \label{eq:crf}
\end{equation}

\noindent\textit{Reference:} VDI~2067 Part~1~\cite{VDI2067}, IRENA~\cite{IRENA2020}.

% -----------------------------------------------
\subsection{Component CAPEX}

\begin{equation}
  C_{\text{CAPEX},i}
  = \hat{Q}_{i} \cdot c^{\text{CAPEX}}_{i} \cdot f_{\text{ann}}
  \quad [\text{EUR}]
  \tag{INV-3}
  \label{eq:inv3}
\end{equation}

where $c^{\text{CAPEX}}_{i}$~[EUR/MW] is the specific investment cost.

% -----------------------------------------------
\subsection{Activation Cost (Binary Investment)}

\begin{equation}
  C_{\text{act},i}
  = y_{i} \cdot c^{\text{act}}_{i} \cdot f_{\text{ann}}
  \quad [\text{EUR}]
  \tag{INV-4}
  \label{eq:inv4}
\end{equation}

\noindent\textit{Purpose:} Captures fixed costs independent of capacity (planning, grid
connection, permits).

% -----------------------------------------------
\subsection{Storage Investment}

\noindent Energy capacity cost:
\begin{equation}
  C_{\text{E},j}
  = \hat{E}_{j} \cdot c^{E}_{j} \cdot f_{\text{ann}}
  \quad [\text{EUR/MWh}]
  \tag{INV-5}
  \label{eq:inv5}
\end{equation}

\noindent Power capacity cost:
\begin{equation}
  C_{\text{P},j}
  = \hat{P}_{j} \cdot c^{P}_{j} \cdot f_{\text{ann}}
  \quad [\text{EUR/MW}]
  \tag{INV-6}
  \label{eq:inv6}
\end{equation}

% -----------------------------------------------
\subsection{Tie-Breaker Cost (Numerical Stability)}

A small non-annualized cost $\epsilon \ll 1$ added per unit of capacity to break degeneracy:

\begin{equation}
  C_{\text{tie},i} = \hat{Q}_{i} \cdot \epsilon_{i}
  \tag{INV-7}
  \label{eq:inv7}
\end{equation}

\noindent\textit{Purpose:} Avoids degenerate solutions where multiple capacity values yield
identical objectives~\cite{Glover1974}.

% =============================================================================
\section{CO\texorpdfstring{$_2$}{2} Emissions Model}
\label{sec:co2}

\textit{Source file:} \texttt{calion/models/emissions\_calculator.py}.

% -----------------------------------------------
\subsection{Grid Electricity Emissions}

\begin{equation}
  \mathrm{CO}_{2,\text{grid},t}
  = P^{\text{el}}_{t} \cdot e^{\text{grid}}_{t} \cdot \Delta t
  \quad [\text{kg}]
  \tag{CO2-1}
  \label{eq:co2-1}
\end{equation}

where $e^{\text{grid}}_{t}$~[kg\,CO$_2$/MWh] is the time-varying grid emission intensity
(marginal or average, depending on methodology).

\noindent\textbf{Annual CO$_2$ cost:}
\begin{equation}
  C_{\text{CO}_2}
  = \frac{\lambda_{\text{CO}_2}}{1000}
    \cdot \sum_{t} \mathrm{CO}_{2,t}
  \quad [\text{EUR}]
  \label{eq:co2-cost}
\end{equation}

where $\lambda_{\text{CO}_2}$~[EUR/tonne\,CO$_2$] is the carbon price.

\noindent\textit{Reference:} For marginal vs.\ average emission accounting, see Tranberg et
al.~\cite{Tranberg2019}, Millar et al.~\cite{Millar2021}.

% -----------------------------------------------
\subsection{Fuel Combustion Emissions}

\begin{equation}
  \mathrm{CO}_{2,\text{fuel},t}
  = F_{t} \cdot e_{f} \cdot \Delta t
  \quad [\text{kg}]
  \tag{CO2-2}
  \label{eq:co2-2}
\end{equation}

where $e_{f}$~[kg\,CO$_2$/MWh$_{\text{fuel}}$] is the specific emission factor of fuel~$f$
(LHV basis).

\noindent\textbf{Standard emission factors (IPCC~2006~\cite{IPCC2006}):}
\begin{itemize}
  \item Natural gas: 201.6~kg\,CO$_2$/MWh ($\approx$ 56~kg\,CO$_2$/GJ)
  \item Light heating oil: 266.4~kg\,CO$_2$/MWh
  \item Biomass (wood chips): 0~kg\,CO$_2$/MWh (carbon neutral, lifecycle considerations aside)
\end{itemize}

% -----------------------------------------------
\subsection{CHP Emission Allocation}

For combined heat and power, total fuel emissions are allocated by output efficiency fractions:

\begin{equation}
  \alpha_{\text{heat}}
  = \frac{\eta^{\text{th}}}{\eta^{\text{th}} + \eta^{\text{el}}},
  \qquad
  \alpha_{\text{el}}
  = \frac{\eta^{\text{el}}}{\eta^{\text{th}} + \eta^{\text{el}}}
  \tag{CO2-3}
  \label{eq:co2-3}
\end{equation}

\begin{align*}
  \mathrm{CO}_{2,\text{heat}} &= \alpha_{\text{heat}} \cdot \mathrm{CO}_{2,\text{total}} \\
  \mathrm{CO}_{2,\text{el}}   &= \alpha_{\text{el}}   \cdot \mathrm{CO}_{2,\text{total}}
\end{align*}

\noindent\textit{Reference:} This efficiency method is one of several allocation methods;
alternatives include the IEA method and bonus method. See Directive~2012/27/EU~\cite{EU2012},
Streckien\.{e} et al.~\cite{Streckiene2009}.

% -----------------------------------------------
\subsection{CO\texorpdfstring{$_2$}{2} Resolution Analysis}

\begin{equation}
  e^{\text{annual}}
  = \frac{\sum_t e^{\text{hourly}}_{t} \cdot \Delta t}{T \cdot \Delta t}
  \quad [\text{kg\,CO}_2/\text{MWh}]
  \label{eq:co2-resolution}
\end{equation}

Resolution scenarios: annual (1 factor), monthly (12), daily (365), hourly (8,760),
15-min (35,040).

\noindent\textit{Reference:} Millar et al.~\cite{Millar2021}, Tranberg et al.~\cite{Tranberg2019}.

% =============================================================================
\section{Sensitivity Analysis}
\label{sec:sensitivity}

\textit{Source file:} \texttt{calion/analysis/sensitivity.py}.

% -----------------------------------------------
\subsection{Parameter Variation Modes}

\noindent Multiplicative variation:
\begin{equation}
  v_k = v_0 \cdot \mu_k
  \tag{SA-1}
  \label{eq:sa1}
\end{equation}

\noindent Absolute variation:
\begin{equation}
  v_k = v_0 + \delta_k
  \tag{SA-2}
  \label{eq:sa2}
\end{equation}

where $v_0$ is the base-case parameter value, $\mu_k$ is the variation multiplier, and
$\delta_k$ is the additive increment.

% -----------------------------------------------
\subsection{Normalized Sensitivity Index}

\begin{equation}
  S_i
  = \frac{|\max_k Z(v_k) - \min_k Z(v_k)|}{Z(v_0)}
  \tag{SA-3}
  \label{eq:sa3}
\end{equation}

where $Z(\cdot)$ is the objective value (total cost) and $v_0$ is the baseline parameter value.

\noindent\textit{Physical meaning:} $S_i$ is the normalized range of objective variation due to
parameter~$i$; analogous to a normalized first-order sensitivity
coefficient~\cite{Saltelli2008}.

\noindent\textit{Reference:} Saltelli et al.~\cite{Saltelli2008} Chapter~1; one-at-a-time (OAT)
sensitivity analysis.

% -----------------------------------------------
\subsection{Standard Sensitivity Parameters}

\begin{table}[h!]
  \centering
  \caption{Standard sensitivity parameters, base values, and variation ranges.}
  \begin{tabular}{@{}llll@{}}
    \toprule
    Parameter & Symbol & Base Value & Range \\
    \midrule
    Gas price               & $\lambda_{\text{gas}}$       & 58.6 EUR/MWh & $\pm$20\,\% \\
    Electricity price       & $\lambda_{\text{el}}$        & 50 EUR/MWh   & $\pm$20\,\% \\
    Heat pump COP factor    & $\eta$                       & 0.75         & $\pm$5\,\%  \\
    P2H efficiency          & $\eta_{\text{P2H}}$          & 0.99         & $\pm$3\,\%  \\
    Storage hourly loss     & $\dot{\lambda}$              & 0.05\,\%/h   & $\pm$50\,\% \\
    Storage charge eff.     & $\eta_c$                     & 0.98         & $\pm$5\,\%  \\
    Storage discharge eff.  & $\eta_d$                     & 0.98         & $\pm$5\,\%  \\
    \bottomrule
  \end{tabular}
  \label{tab:sensitivity}
\end{table}

% =============================================================================
\section{Notation Summary}
\label{sec:notation}

\begin{table}[h!]
  \centering
  \caption{Notation summary.}
  \begin{tabular}{@{}lll@{}}
    \toprule
    Symbol & Unit & Description \\
    \midrule
    $T$                              & --             & Number of timesteps \\
    $\Delta t$                       & h              & Timestep duration \\
    $Q_t$                            & MW             & Thermal power \\
    $P^{\text{el}}_t$                & MW             & Electrical power \\
    $F_t$                            & MW             & Fuel input power \\
    $E_t$                            & MWh            & Stored energy \\
    $\hat{Q}$                        & MW             & Installed thermal capacity \\
    $\hat{E}$                        & MWh            & Installed energy capacity \\
    $\hat{P}$                        & MW             & Installed power capacity \\
    $u_t, \xi_t$                     & --             & Binary operating mode indicator \\
    $\xi^{\text{active}}_t$          & --             & Storage activation binary \\
    $Q^{\text{slack}}_t$             & MW             & Demand slack at consumer node (penalized) \\
    $\lambda_t$                      & --             & Normalized load fraction $Q_t / Q^{\text{peak}}$ \\
    $\lambda^{(k)}, T^{(k)}_s, T^{(k)}_r$ & --, $^\circ$C & PWL breakpoint $k$ (load fraction, temperatures) \\
    $y$                              & --             & Binary investment decision \\
    $\COP_t$                         & --             & Coefficient of performance \\
    $\eta^{\text{th}}, \eta^{\text{el}}$ & --         & Thermal / electrical efficiency \\
    $\eta_c, \eta_d$                 & --             & Charge / discharge efficiency \\
    $\lambda$                        & --             & Self-discharge factor per hour \\
    $L$                              & m              & Pipe length \\
    $U$                              & W/(m$\cdot$K)  & Linear heat transfer coefficient \\
    $\dot{m}$                        & kg/s           & Mass flow rate \\
    $c_p$                            & kJ/(kg$\cdot$K)& Specific heat capacity (water: 4.186) \\
    $T_s, T_r, T_g$                  & $^\circ$C      & Supply / return / ground temperature \\
    $e_f$                            & kg\,CO$_2$/MWh & Fuel CO$_2$ emission factor \\
    $e^{\text{grid}}_t$              & kg\,CO$_2$/MWh & Grid CO$_2$ intensity at timestep $t$ \\
    $\lambda_{\text{CO}_2}$          & EUR/t\,CO$_2$  & Carbon price \\
    $c^{\text{CAPEX}}$               & EUR/MW         & Specific investment cost \\
    $\tau_{\text{life}}$             & yr             & Economic asset lifetime \\
    $f_{\text{ann}}$                 & --             & Annualization factor \\
    $M$                              & --             & Big-M constant (large number) \\
    \bottomrule
  \end{tabular}
  \label{tab:notation}
\end{table}

% =============================================================================
\section*{Scientific References}
\addcontentsline{toc}{section}{Scientific References}

\begin{thebibliography}{99}

% MILP Optimization and Energy Systems
\bibitem{Morales2014}
Morales, J.M., Conejo, A.J., Madsen, H., Pinson, P., Zugno, M.\ (2014).
\textit{Integrating Renewables in Electricity Markets: Operational Problems}.
Springer, New York. ISBN 978-1-4614-9411-9.

\bibitem{Savelsbergh1994}
Savelsbergh, M.\ (1994).
Preprocessing and probing techniques for mixed integer programming problems.
\textit{ORSA Journal on Computing}, 6(4), 445--454.
\url{https://doi.org/10.1287/ijoc.6.4.445}

\bibitem{Lund2014}
Lund, H., Werner, S., Wiltshire, R., Svendsen, S., Thorsen, J.E., Hvelplund, F.,
Mathiesen, B.V.\ (2014).
4th Generation District Heating (4GDH): Integrating smart thermal grids into future
sustainable energy systems.
\textit{Energy}, 68, 1--11.
\url{https://doi.org/10.1016/j.energy.2014.02.089}

\bibitem{Connolly2014}
Connolly, D., Lund, H., Mathiesen, B.V., Werner, S., M\"{o}ller, B., Persson, U.,
Boermans, T., Trier, D., \O{}stergaard, P.A., Nielsen, S.\ (2014).
Heat Roadmap Europe: Combining district heating with heat savings to decarbonise the EU
energy system.
\textit{Energy Policy}, 65, 475--489.
\url{https://doi.org/10.1016/j.enpol.2013.10.035}

\bibitem{Henning2014}
Henning, H.M., Palzer, A.\ (2014).
A comprehensive model for the German electricity and heat sector in a future energy system
with a dominating contribution from renewable energy technologies --- Part I: Methodology.
\textit{Renewable and Sustainable Energy Reviews}, 30, 1003--1018.
\url{https://doi.org/10.1016/j.rser.2013.09.012}

% Electricity Market and Grid Pricing
\bibitem{Kirschen2018}
Kirschen, D., Strbac, G.\ (2018).
\textit{Fundamentals of Power System Economics} (2nd ed.). Wiley-Blackwell.
ISBN 978-1-119-02035-1.

\bibitem{Pfenninger2014}
Pfenninger, S., Hawkes, A., Keirstead, J.\ (2014).
Energy systems modeling for twenty-first century energy challenges.
\textit{Renewable and Sustainable Energy Reviews}, 33, 74--86.
\url{https://doi.org/10.1016/j.rser.2014.02.003}

\bibitem{BNetzA2023}
Bundesnetzagentur\ (2023).
\textit{Netzentgeltsystematik Strom} (Grid tariff methodology).
Bonn: Federal Network Agency Germany. \url{www.bundesnetzagentur.de}

% Heat Pump Technology and COP
\bibitem{Vesterlund2017}
Vesterlund, M., Toffolo, A., Dahl, J.\ (2017).
Optimization of multi-source complex district heating network, a case study.
\textit{Energy}, 126, 53--63.
\url{https://doi.org/10.1016/j.energy.2017.02.148}

\bibitem{Lund2016b}
Lund, R., Persson, U.\ (2016).
Mapping of potential heat sources for heat pumps for district heating in Denmark.
\textit{Energy}, 110, 129--138.
\url{https://doi.org/10.1016/j.energy.2015.12.127}

\bibitem{Viessmann2021}
Viessmann\ (2021).
\textit{Heat Pump Technology --- Planning and Specification Guide}.
Viessmann Climate Solutions SE, Allendorf (Eder).

\bibitem{Berghmans2012}
Berghmans, J.\ (2012).
Heat pumps. In: \textit{Energy Conversion} (A.\ De Vos, ed.). Wiley-VCH.

\bibitem{Wolf2014}
Wolf, S., Fahl, U., Blesl, M., Vo\ss{}, A., Jakobs, R.\ (2014).
Analyse des Potenzials von Industriew\"{a}rmepumpen in Deutschland.
\textit{IER Forschungsbericht}, 113. Universit\"{a}t Stuttgart.

\bibitem{Morvaj2016}
Morvaj, B., Evins, R., Carmeliet, J.\ (2016).
Optimising urban energy systems: Simultaneous system sizing, operation and district
heating network layout optimization.
\textit{Energy}, 116(1), 619--636.
\url{https://doi.org/10.1016/j.energy.2016.09.139}

\bibitem{Raman1994}
Raman, R., Grossmann, I.E.\ (1994).
Modelling and computational techniques for logic based integer programming.
\textit{Computers \& Chemical Engineering}, 18(7), 563--578.
\url{https://doi.org/10.1016/0098-1354(93)E0010-7}

\bibitem{Staffell2012}
Staffell, I., Brett, D., Brandon, N., Hawkes, A.\ (2012).
A review of domestic heat pumps.
\textit{Energy \& Environmental Science}, 5(11), 9291--9306.
\url{https://doi.org/10.1039/C2EE22653G}

\bibitem{Ruhnau2019}
Ruhnau, O., Hirth, L., Praktiknjo, A.\ (2019).
Time series of heat demand and heat pump efficiency for energy system modeling.
\textit{Scientific Data}, 6, 189.
\url{https://doi.org/10.1038/s41597-019-0199-y}

\bibitem{Cengel2018}
\c{C}engel, Y.A., Boles, M.A.\ (2018).
\textit{Thermodynamics: An Engineering Approach} (9th ed.). McGraw-Hill.
ISBN 978-0-07-339817-4.

\bibitem{Incropera2017}
Incropera, F.P., DeWitt, D.P., Bergman, T.L., Lavine, A.S.\ (2017).
\textit{Fundamentals of Heat and Mass Transfer} (8th ed.). Wiley.
ISBN 978-0-471-45728-2.

\bibitem{EN14825}
DIN EN~14825:2022 --- Air conditioners, liquid chilling packages and heat pumps, with
electrically driven compressors, for space heating and cooling --- Testing and rating at
part load conditions and calculation of seasonal performance.
Beuth Verlag, Berlin.

% Thermal Storage
\bibitem{Quoilin2017}
Quoilin, S., Hidalgo Gonzalez, I., Zucker, A.\ (2017).
\textit{Modelling Future EU Power Systems Under High Shares of Renewables}.
JRC Technical Report EUR~28380~EN. Publications Office of the EU.

\bibitem{Lund2007}
Lund, H., Duic, N., Krajacic, G., da Gra\c{c}a Carvalho, M.\ (2007).
Two energy system analysis models: A comparison of methodologies and results.
\textit{Energy}, 32(6), 948--954.
\url{https://doi.org/10.1016/j.energy.2006.08.014}

\bibitem{Haller2019}
Haller, M.Y., Damerau, K., Dreier, D., Marty, H., Herkel, S., Jenni, A.\ (2019).
Comparison of Control Strategies for District Thermal Energy Storage.
\textit{Energies}, 12(16), 3060.
\url{https://doi.org/10.3390/en12163060}

% Power-to-Heat
\bibitem{Lund2015}
Lund, P.D., Lindgren, J., Mikkola, J., Salpakari, J.\ (2015).
Review of energy system flexibility measures to enable high levels of variable renewable
electricity.
\textit{Renewable and Sustainable Energy Reviews}, 45, 785--807.
\url{https://doi.org/10.1016/j.rser.2015.01.057}

% Thermal Generators and CHP
\bibitem{Eurostat2023}
Eurostat\ (2023).
\textit{Energy Efficiency --- Thermal Plant Efficiencies}.
European Commission Statistics. \url{https://ec.europa.eu/eurostat}

% Pipe Heat Loss and District Heating
\bibitem{EN13941}
DIN EN~13941-1:2019 --- District heating pipes --- Design and installation.
Beuth Verlag, Berlin.

\bibitem{Frederiksen2013}
Frederiksen, S., Werner, S.\ (2013).
\textit{District Heating and Cooling}. Studentlitteratur, Lund.
ISBN 978-91-44-08530-2.

\bibitem{Benonysson1995}
Benonysson, A., B\o{}hm, B., Ravn, H.F.\ (1995).
Operational optimization in a district heating system.
\textit{Energy Conversion and Management}, 36(5), 297--314.
\url{https://doi.org/10.1016/0196-8904(95)98895-T}

\bibitem{Mesfun2015}
Mesfun, S., Toffolo, A.\ (2015).
Optimization of a complex district heating system using a MILP model.
\textit{ECOS 2015 --- 28th International Conference on Efficiency, Cost, Optimization,
Simulation and Environmental Impact of Energy Systems}, Pau, France.

\bibitem{VDI2067}
VDI~2067 Part~1:2012 --- Economic efficiency of building installations --- Fundamentals
and economic calculation.
VDI-Verlag, D\"{u}sseldorf.

\bibitem{Williams2013}
Williams, H.P.\ (2013).
\textit{Model Building in Mathematical Programming} (5th ed.). Wiley.
ISBN 978-1-118-44333-0.

\bibitem{Beale1970}
Beale, E.M.L., Tomlin, J.A.\ (1970).
Special facilities in a general mathematical programming system for non-convex problems
using ordered sets of variables.
\textit{5th International Conference on Operations Research}, J.\ Lawrence (ed.),
Tavistock, London, 447--454.

\bibitem{Schweiger2017}
Schweiger, G., Larsson, P.O., Magnusson, F., Lauenburg, P., Velut, S.\ (2017).
District heating and cooling systems --- Framework for Modelica-based simulation and
dynamic optimization.
\textit{Energy}, 137, 566--578.
\url{https://doi.org/10.1016/j.energy.2017.05.115}

% Electricity Market (continued)
\bibitem{Hirth2013}
Hirth, L.\ (2013).
The market value of variable renewables.
\textit{Energy Economics}, 38, 218--236.
\url{https://doi.org/10.1016/j.eneco.2013.02.004}

% Investment
\bibitem{IRENA2020}
IRENA\ (2020).
\textit{Renewable Power Generation Costs in 2019}.
International Renewable Energy Agency, Abu Dhabi.
ISBN 978-92-9260-040-2.

\bibitem{Glover1974}
Glover, F., Woolsey, E.\ (1974).
Converting the 0-1 polynomial programming problem to a 0-1 linear program.
\textit{Operations Research}, 22(1), 180--182.
\url{https://doi.org/10.1287/opre.22.1.180}

% CO2 Emissions
\bibitem{Tranberg2019}
Tranberg, B., Corradi, O., Lajoie, B., Gibon, T., Staffell, I., Andresen, G.B.\ (2019).
Real-time carbon accounting method for the European electricity markets.
\textit{Energy Strategy Reviews}, 26, 100367.
\url{https://doi.org/10.1016/j.esr.2019.100367}

\bibitem{Millar2021}
Millar, M.A., Yu, Z., Burnside, N., Jones, G., Elrick, B.\ (2021).
Identification of key performance indicators and complimentary load profiles for the
integration of low-carbon technologies in existing district networks.
\textit{Applied Energy}, 282, 116108.
\url{https://doi.org/10.1016/j.apenergy.2020.116108}

\bibitem{IPCC2006}
IPCC\ (2006).
\textit{2006 IPCC Guidelines for National Greenhouse Gas Inventories, Volume~2: Energy}.
Eggleston H.S., Buendia L., Miwa K., Ngara T., Tanabe K.\ (eds). IGES, Japan.

\bibitem{EU2012}
European Parliament and Council\ (2012).
\textit{Directive 2012/27/EU on energy efficiency}.
Official Journal of the European Union, L~315/1.

\bibitem{Streckiene2009}
Streckien\.{e}, G., Martinaitis, V., Andersen, A.N., Katz, J.\ (2009).
Feasibility of CHP-plants with thermal stores in the German spot market.
\textit{Applied Energy}, 86(12), 2308--2316.
\url{https://doi.org/10.1016/j.apenergy.2009.03.011}

% Sensitivity Analysis
\bibitem{Saltelli2008}
Saltelli, A., Ratto, M., Andres, T., Campolongo, F., Cariboni, J., Gatelli, D.,
Saisana, M., Tarantola, S.\ (2008).
\textit{Global Sensitivity Analysis: The Primer}. Wiley.
ISBN 978-0-470-05997-5.

\end{thebibliography}

\end{document}
